% Abstract: at most 200 words. No number that is not in the body. Every number has a % src.
\begin{abstract}
AlphaFold2 is written in JAX. JAX treats each chip of a tensor processing unit (TPU) slice as a separate device. A JAX
program does not spread its work over the chips by itself, so a user who runs AlphaFold2 on a slice can pay for chips
that do nothing. We found no published measurement of how AlphaFold2 inference uses a multi-chip TPU slice. Our goal
was to measure how much of an eight-chip TPU v5e slice AlphaFold2 inference uses, and which way of mapping the work in
JAX uses all of it. Across 12 experiments, the model used random weights, so we studied execution, not prediction
quality. % src: recomputed: COUNT.tpu_experiments_with_raw_data
By default, only one chip held data.
Batching up to eight input sequences (queries) of 118 residues on one chip with \texttt{jax.vmap} did not raise throughput. % src: results/sweep/batching_sweep.json -> [*].batch_size; recomputed: INPUT.toy_length
Mapping independent queries to the chips with \texttt{jax.pmap} ran them in parallel on all eight chips.
On a grid of 1 to 8 chips and 100 to 1000 residues, eight chips gave 6.5 to 7.9 times the throughput of one chip. % src: results/sweep/scaling_law_data.json -> grid[*].chips, grid[*].length; recomputed: V2.speedup_8over1_min, V2.speedup_8over1_max
JAX's automatic partitioning did not shorten the call for one query. Our program did not tell the partitioner how to
split any array across the chips.
\end{abstract}

\begin{IEEEkeywords}
AlphaFold2, tensor processing units, JAX, data parallelism, protein structure prediction
\end{IEEEkeywords}
